%% =====================================================================
%%  V2 DRAFT SECTIONS — reframed around the defensible novel core.
%%  These are drop-in replacements / additions for article-v1.tex.
%%  Every number here is reproduced by the evaluation/ and improvement/
%%  modules (see results/*.md). Rationale for each change is in the
%%  comment above it. NOT auto-merged — review and paste per section.
%%
%%  Summary of what changed from v1 and WHY:
%%   1. Abstract + contributions: the novelty is no longer "a novel
%%      rule-based approach" (the audit found rule-based hybrids, WBF, and
%%      appearance filtering are all established). The novelty is the
%%      QUANTIFIED complementarity: traditional cues recover trees beyond
%%      the deep-ensemble detection ceiling, characterised mechanistically.
%%   2. Results: the recall-only "96%% / +4.8%%" headline is replaced by a
%%      full precision/recall/F1 + mAP table under one-to-one matching. The
%%      original 96%% is retained honestly as the v1 metric it reproduces.
%%   3. Honest trade-off + the precision gate that resolves it.
%%   4. New limitations subsection.
%% =====================================================================


%% ---------------------------------------------------------------------
%% REPLACEMENT ABSTRACT
%% Why: v1 abstract sells "a novel rule-based approach" and frames the
%% result as "increase the number of detected tree crowns" (recall only).
%% v2 leads with the quantified beyond-ceiling finding and reports both
%% recall AND precision so the contribution survives review.
%% ---------------------------------------------------------------------
\begin{abstract}
Global warming, loss of biodiversity, and air pollution are among the most
significant problems facing Earth, and monitoring forests is central to
addressing them. Automatic tree crown detection from satellite imagery is
therefore an important task, approached by both traditional image-processing
and deep-learning methods. In this study we ask a specific question: once a
strong ensemble of deep detectors has been applied, is there anything a
traditional method can still contribute? We combine five deep detectors with
weighted boxes fusion (WBF) and a traditional pipeline (illumination-invariant
colour and a Gabor filter bank feeding a watershed/local-maxima detector)
through a rule-based integration step. Evaluating on a benchmark of 222 test
images containing 13{,}552 trees under one-to-one matching, we report precision,
recall, and $F_1$ for every method rather than recall alone. Our central
finding is that the traditional step recovers 448 true crowns (3.3 percentage
points) that lie \emph{beyond the detection ceiling of the full six-model
deep-learning stack at its most permissive thresholds} --- trees no detector
proposes at any confidence. We show these recovered crowns are systematically
small (a median crown area $2.44\times$ below detected crowns) and spectrally
faint, which explains why appearance-trained detectors miss them and a
texture/local-maxima operator recovers them. We further show that a lightweight
appearance gate on the traditional-only detections removes the non-tree
false positives they introduce, so the integrated detector dominates the WBF
ensemble on the precision--recall frontier ($F_1$ 0.921 $\rightarrow$ 0.934).
We report the advantages, the precision cost, and the limitations of the
approach.
\end{abstract}


%% ---------------------------------------------------------------------
%% ADD to the end of the Introduction: an explicit CONTRIBUTIONS list.
%% Why: v1 never states contributions explicitly, and calls the method
%% "novel" without isolating what is actually new vs. established. This
%% list is honest about which parts are contributions and which are
%% competent use of known tools.
%% ---------------------------------------------------------------------
\paragraph{Contributions.} We make the following contributions, and are
explicit about their novelty:
\begin{itemize}
  \item \textbf{(Primary, novel)} A quantified, benchmark-backed measurement of
  the \emph{complementarity} between traditional and deep detectors for tree
  crowns: a traditional operator recovers 3.3\% of trees beyond the full
  deep-ensemble ceiling, and we characterise \emph{why} (small, spectrally
  faint crowns). Prior hybrid studies assert complementarity qualitatively;
  we measure it.
  \item \textbf{(Methodological)} A rule-based integration procedure that
  decides, per candidate, whether to trust the detector, a neighbour, or the
  traditional evidence. The individual operations are standard; the codified
  decision procedure and its evaluation are the contribution.
  \item \textbf{(Evaluation)} A rigorous re-evaluation protocol
  (one-to-one matching with precision/recall/$F_1$ and COCO mAP, plus a
  precision--recall operating-curve analysis) that corrects the recall-only
  scoring common in this sub-area.
  \item \textbf{(Engineering)} A faithful implementation of the traditional
  features the method requires (a multi-orientation, multi-scale Gabor bank and
  an illumination-invariant colour feature) and an appearance gate that removes
  the false positives introduced by the traditional step.
\end{itemize}


%% ---------------------------------------------------------------------
%% REPLACEMENT for Section 5.3 "Performance of Deep Learning Methods"
%% Why: v1 reports only "12353 / 91.2%%" with no precision. v2 gives the
%% full per-method table under one-to-one matching + standard mAP.
%% ---------------------------------------------------------------------
\subsection{Performance of Deep Learning Methods}

We evaluate each detector and the WBF ensemble under one-to-one (Hungarian)
matching between predictions and ground-truth boxes at an intersection-over-union
threshold of 0.5, which --- unlike a recall-only count --- penalises false
positives. We report precision, recall, and $F_1$ at the operating point used by
the integration pipeline (detector score $\geq 0.8$), and COCO mean average
precision (AP, AP$_{50}$) computed over all detections (threshold-free).
Table~\ref{tab:results} summarises the results on the 13{,}552-tree test set.

\begin{table}[htbp]
\centering
\caption{Detection performance on the test set (222 images, 13{,}552 trees).
Precision/recall/$F_1$ are one-to-one matched at IoU $0.5$, detector score
$\geq 0.8$; AP/AP$_{50}$ are standard COCO mAP over all detections. The
integrated method outputs points, so mAP is not applicable.}
\label{tab:results}
\begin{tabular}{lccccc}
\toprule
Method & Precision & Recall & $F_1$ & AP & AP$_{50}$ \\
\midrule
Swin Transformer   & 0.979 & 0.810 & 0.886 & 0.504 & 0.856 \\
Cascade R-CNN      & 0.984 & 0.786 & 0.874 & 0.512 & 0.855 \\
Deformable DETR    & \textbf{0.992} & 0.725 & 0.838 & 0.492 & 0.856 \\
Faster R-CNN       & 0.974 & 0.822 & 0.892 & 0.500 & 0.846 \\
YOLOv3             & 0.985 & 0.549 & 0.705 & 0.477 & 0.848 \\
WBF ensemble       & 0.968 & 0.878 & 0.921 & 0.520 & \textbf{0.864} \\
\midrule
Integrated (this work)        & 0.833 & 0.971 & 0.897 & n/a & n/a \\
Integrated + appearance gate  & 0.962 & 0.908 & \textbf{0.934} & n/a & n/a \\
\bottomrule
\end{tabular}
\end{table}

Among individual detectors, Faster R-CNN attains the best $F_1$ (0.892). The WBF
ensemble improves over every single model ($F_1$ 0.921, recall 0.878 at
precision 0.968), confirming the value of multi-model fusion reported in our
prior work~\cite{durgut2025multi}. Notably, all detectors saturate near
AP$_{50}\approx 0.85$; the ensemble's advantage is in recall at high precision,
not in a higher AP$_{50}$.

For reference and reproducibility, under the original recall-only protocol used
in the conference version of this work (a ground-truth crown is counted as
detected if any prediction centre falls within it), the ensemble reaches
91.2\% (12{,}353/13{,}552) and the integrated method 96.0\% (13{,}012/13{,}552).
We report the one-to-one metrics above as the primary results because the
recall-only protocol does not account for false positives.


%% ---------------------------------------------------------------------
%% REPLACEMENT for Section 5.4 "Performance of Rule-based Algorithm"
%% Why: v1 attributes the gain to "solving small/close trees" from figures
%% only. v2 proves the beyond-ceiling recovery quantitatively, gives the
%% mechanism, and is honest about the precision cost + the gate that fixes it.
%% ---------------------------------------------------------------------
\subsection{Performance of the Rule-based Integration}

\paragraph{Recovery beyond the deep-ensemble ceiling.}
The key question is whether the traditional step adds detections that the deep
ensemble fundamentally cannot produce, as opposed to detections that a lower
score threshold would also yield. To test this, we take the union of all six
detectors' boxes \emph{at their most permissive score thresholds} --- the
theoretical ceiling of what the deep-learning stack can propose --- and measure
its ground-truth coverage. This union reaches 94.2\% of trees; the WBF ensemble
reaches 92.5\%; the best single detector reaches at most 86.9\%. The integrated
method reaches 97.5\%. Thus the traditional step recovers a net 448 crowns
(3.3\%) that \emph{no} detector proposes at \emph{any} confidence, and a
beyond-ceiling set of 692 crowns hit by the integration but by no detector. The
common objection ``simply lower the detector threshold'' is therefore
insufficient: every detector is already at its floor.

\paragraph{Why the detectors miss these crowns.}
Characterising the recovered crowns against the detected population reveals a
consistent pattern (Table~\ref{tab:beyond}). Beyond-ceiling crowns have a median
area of 788\,px$^2$ versus 1924\,px$^2$ for detected crowns --- $2.44\times$
smaller, with 82\% falling below the 25th percentile of detected-crown area ---
and are brighter, less saturated, and lower in excess-green, i.e.\ small and
spectrally faint. They retain higher local texture contrast, which is precisely
the cue the Gabor/local-maxima operator responds to. Appearance-trained
detectors are least confident on small faint crowns and drop them below any
usable threshold, whereas a texture/extremum operator, which does not rely on
learned appearance priors, still fires. A visual sample confirms the recovered
detections are genuine small crowns rather than annotation artefacts.

\begin{table}[htbp]
\centering
\caption{Median properties of ground-truth crowns reachable by the deep-detector
union versus those recovered only beyond that ceiling.}
\label{tab:beyond}
\begin{tabular}{lcc}
\toprule
 & DL-reachable (12{,}762) & Beyond-ceiling (692) \\
\midrule
Crown area (px$^2$)     & 1924 & 788 \\
Brightness (HSV $V$)    & 95.1 & 102.5 \\
Saturation              & 86.0 & 59.9 \\
Excess-green            & 22.3 & 11.1 \\
Local texture (std)     & 26.8 & 32.4 \\
\bottomrule
\end{tabular}
\end{table}

\paragraph{The precision cost and how we resolve it.}
Recovering faint crowns with a traditional operator also introduces false
positives: the traditional-only detections fire on grass, bare soil, and
rooftops, so the integrated method trades precision (0.833) for recall (0.971)
relative to the ensemble (Table~\ref{tab:results}). Of its 2{,}647 false
positives, 85\% are traditional-only points far from any reliable detector box.
We therefore add a lightweight appearance gate: for each traditional-only
candidate we classify a crown-sized crop as tree vs.\ non-tree from
brightness, chromaticity, and texture features, keeping detections near a
reliable detector box unconditionally. Trained and evaluated on
image-disjoint halves of the test set, the gate raises precision to 0.962 while
retaining recall 0.908, giving $F_1 = 0.934$ --- above the WBF ensemble's 0.921
at a recall the ensemble cannot reach. The integrated detector thus
\emph{dominates} the ensemble on the precision--recall frontier rather than
merely trading along it.


%% ---------------------------------------------------------------------
%% NEW subsection to add before Conclusion.
%% Why: v1 has no limitations section; reviewers expect one, and stating
%% these ourselves is far stronger than having them raised against us.
%% ---------------------------------------------------------------------
\subsection{Limitations}
Three limitations bound our claims. First, the complementarity is quantified on
a single dataset covering the Marmara and Mediterranean regions of Turkey; the
$2.44\times$-smaller mechanism is plausibly general, but the specific 3.3\%
magnitude and the appearance-gate coefficients are dataset-dependent and require
re-fitting for other biomes. Second, ``beyond-ceiling'' is defined relative to
these six detectors; a detector trained specifically for tiny crowns could
shrink the gap, so the finding characterises the complementarity of the current
strong ensemble rather than a permanent limit of deep detection. Third, the
appearance gate is supervised; for unlabelled imagery it must be trained once on
a labelled subset, which the image-disjoint evaluation here is designed to
estimate.
